\subsection{Transfer of Pieri expansions to all nonzero-Gram cells}
\label{sec:all-cell-pieri-transfer}

We retain the original Bessel subgroups, characters, compact
types, quotient gauges and radial Jacquet norm of the signed
Pieri-family and iterated-minor proofs in the construction
record~\cite{current-construction-record}.
The new step is that a
coefficient with zero radial Jacquet norm has zero orbital
integral on every relevant nonzero-Gram cell, not just on
the open cell. This permits a direct weighted-sum argument.

\paragraph{A two-sided relative functional.}
Write \(S=K_zN\), \(S'=K_{z'}N\), and
\(\tau(wn)=\pi_W(w)\psi(n)\), with the analogous \(\tau'\).
These are unitary representations of the indicated
semidirect products on the compact-type spaces \(W,W'\):
the compact factors normalize the flag unipotents and
preserve their characters. Their Haar measures are the
original compact and unipotent product measures. Use
inner products linear in the first variable.

Let \(\sigma\) be a discrete series with the original
large-gap decay on smooth vectors. Its two original
period-normalized maps are denoted \(B:H^\infty\to W\),
\(B':H^\infty\to W'\). In particular, for the first map
the established period factorization is
\[
 \int_S\langle\sigma(s)x,v\rangle
                \overline{\langle\tau(s)a,u\rangle}\,ds
       =\langle Bx,a\rangle\overline{\langle Bv,u\rangle}.
                                                        \tag{AP1}
\]
The formula includes the case \(B=0\); likewise for
\(S'\). We use the multiplicity-one theorem for these
smooth Bessel functionals and the mixed Jacquet identity
already used in the local Pieri family and Jacquet norm
proofs. In particular a coefficient
\[
 F(h)=\sum_j\langle\sigma(h)y_j,x_j\rangle A_j
\]
has original squared radial norm
\[
 jNorm2(F)=c_{\rm rad}d_\sigma^{-1}|s_\sigma(F)|^2,\qquad
 s_\sigma(F)=\sum_j\langle A_jB'y_j,Bx_j\rangle,          \tag{AP2}
\]
with the same positive radial factor \(c_{\rm rad}\)
and formal degree \(d_\sigma\) as before.

Fix any represented relevant point with nonzero Gram
and representative \(g\). Put \(R=N\cap gN g^{-1}\),
\(R'=g^{-1}Rg\). Give \(R\) exactly its prescribed
root-coordinate Haar measure and \(R'\) the transported
measure. The quotient in \(\mathrm{OrbInt}\) is
\[
 \Delta_g R\backslash(S\times S'),\qquad
       \Delta_g R=\{(a,g^{-1}ag):a\in R\}.
\]
The integrand is
\[
 \tau(s)F(s^{-1}gs')\tau'(s')^{-1}.
\]
It is invariant under the indicated left diagonal action
by relevance. The identification holds for every admitted
\(\mathrm{IsOrbQuot}\) measure: the Weil identity first
forces its support, after pushforward to the orbit
space, to lie on this quotient, and then determines its
pushforward there. More explicitly, test the Weil identity
on compactly supported functions of a quotient coordinate
and a fiber coordinate with integral one. Local sections
and a countable partition of unity prove the assertion
without any requirement that the given measure itself
be supported on a particular section.
The nonzero-Gram root-coordinate chart of the flag
stabilizer is bijective with triangular unipotent
translation law; product Lebesgue measure, with the
specified root scalar, is therefore Haar.

All the integrals below converge absolutely, also for
arbitrary smooth \(x_j,y_j\) in the coefficients, and are
continuous for finite smooth seminorms in these vectors.
Here is the quantitative reason for the use of the
large-gap hypothesis. The terminal quotient properness
estimate in \texttt{global-norm-frontier-finite-decay.tex}
of that record applies
to the two paired flags: after removing the free
stabilizer roots, the real noncompact coordinates \(y\)
have dimension at most \(2n^2\) and
\[
                 1+|y|\le C_g'(1+\|h\|)^{64s^2},
       \qquad n=2r+k+1,\quad s=r+1 ,
\]
where \(h\) is their product with \(g\). This estimate is
for the actual terminal fibre, with its two compact
factors and root measure; it uses the nondegenerate
paired-frame Gram matrix and positive orthogonal
complement, not an inverse of a possibly vanishing
internal minor. The signed block Cramer formulas
used in that estimate apply with any fixed exact
representative \(g\) of the given nonzero-Gram cell;
only the constant \(C'_g\) changes. A coefficient
majorized by \(C(1+\|h\|)^{-A}\) is therefore integrable
if \(A>64s^2(2n^2+1)\). The smooth coefficient bound
in \texttt{global-norm-frontier-pieri-family.tex},
in the original reserve of the shifted family, gives
\(A\ge (C_g/2)s^4(3sn^2+n^3)-1\). Taking
\(C_g\ge10^6\) is more than sufficient. The constants
in this qualitative integrability argument never enter
a uniform norm estimate.

\begin{theorem}[All-cell factorization of coefficient functionals]
\label{thm:all-cell-relative-factorization}
Under these actual period, smooth-decay and quotient
hypotheses, for a fixed relevant nonzero-Gram orbital
datum there is a linear operator \(T_{g,\nu}:W'\to W\)
such that every finite smooth coefficient of \(\sigma\)
satisfies
\[
             \mathrm{OrbInt}(F,g,\nu)
                     =s_\sigma(F)T_{g,\nu}.             \tag{AP3}
\]
In particular \(jNorm2(F)=0\) implies vanishing of
the original orbital integral at every such datum.
No nonvanishing assertion for the full space of Bessel
functionals is needed if a period map is zero.
\end{theorem}

\begin{proof}
Choose a measurable section for \(R'\backslash S'\).
The diagonal quotient disintegrates into the full
first \(S\) and this second quotient. Indeed apply
the two successive Weil identities to a compactly
supported test function on \(S\times S'\): their
composition is exactly the diagonal Weil identity
with the same \(R\) measure. Thus the product measures
on these coordinates have the required normalization.
The reverse disintegration also holds. Both orders
are permitted by absolute convergence.

It suffices to consider
\(F(h)=\langle\sigma(h)y,x\rangle (a\otimes a'^*)\),
where \((a\otimes a'^*)w=\langle w,a'\rangle a\).
Integration over the full first \(S\) uses the complex
conjugate of (AP1) and gives, at \(w\in W'\),
\[
  \langle a,Bx\rangle
        \int_{R'\backslash S'}
           B(\sigma(gs')y)\,
                     \langle w,\tau'(s')a'\rangle\,d\dot s'.
                                                        \tag{AP4}
\]
The integrand is independent of the representative
of \(s'\), since \(B\) is equivariant and relevance
identifies the two characters on \(R,R'\).
If \(B=0\), (AP4) gives zero. If \(B'=0\), integrate the
full second \(S'\) first, using the corresponding
period identity on that side, to obtain zero as well.
This reverse integration is essential: a zero period
form alone has not been identified with the absence
of other equivariant smooth maps.

Suppose now that both maps are nonzero. Fix \(x,a\)
with \(\langle a,Bx\rangle\ne0\). Then the integral
in (AP4) is absolutely convergent and continuous
in \(y\) for the smooth topology, by the bound for
the full orbital integral and this fixed nonzero
factor. For fixed \(w\in W',u\in W\), define a smooth
linear map \(B_{w,u}:H^\infty\to W'\) by
\[
 \langle B_{w,u}y,a'\rangle
   =\int_{R'\backslash S'}
       \langle B(\sigma(gs')y),u\rangle
                    \langle w,\tau'(s')a'\rangle\,d\dot s'.
\]
The right side is linear in \(y\), conjugate linear
in \(a'\); finite dimensionality of \(W'\) gives this
map and its continuity. The quotient measure is
right \(S'\)-invariant, since \(R'\) and \(S'\) are
unimodular. Substitute \(s'\mapsto s't\) to get
\(B_{w,u}\sigma(t)=\tau'(t)B_{w,u}\).
By Bessel multiplicity one,
\(B_{w,u}=c_{w,u}B'\). The scalars \(c_{w,u}\) are
linear in \(w\), conjugate linear in \(u\), hence
are the matrix entries of one operator \(T_{g,\nu}\)
independent of \(x,y,a,a'\).
Substitution in (AP4) gives
\[
 \langle a,Bx\rangle\langle B'y,a'\rangle T_{g,\nu}
       =\langle (a\otimes a'^*)B'y,Bx\rangle T_{g,\nu}.
\]
Summing proves (AP3), and (AP2) proves its last
assertion. The formula is for the same orbital
measure on both sides.
\end{proof}

\paragraph{The same Pieri coefficients on every cell.}
Use the original signed exterior-Pieri decomposition
and its positive leading-minor norm and coefficient-sum
estimates. Fix the finite shift ball allowed by
\(\mathrm{BesselRangeNearParameter}\).
For each parameter with nonzero Jacquet image choose
the original normalized coefficient \(F_\eta\), with
\(jNorm2(F_\eta)=1\); use zero on a zero image.
The established mixed Jacquet identity makes the
nonzero images at distinct parameters orthogonal.
There is no chamber permutation inside this ball,
by the retained gap reserve.

For the leading minor \(D_\ell\), multiplying a
coefficient by \(D_\ell\) is the same finite exterior
tensor coefficient used in the preceding proof.
The minor is constant on the two-sided orbit, since
the frame actions are upper unitriangular and the
compact groups fix the marked frames. Thus
\[
 \mathrm{OrbInt}(D_\ell F_\eta,g,\nu)
            =D_\ell(\gamma)\mathrm{OrbInt}(F_\eta,g,\nu).
\]
The tensor decomposition gives the actual coefficient
in each shifted representation. On the big cell its
Jacquet image agrees with \(a_{\eta,I}^{(\ell)}F_{\eta+\mathbf1_I}\)
after Jacquet integration, by the proved expansion.
Their difference is a coefficient of the same shifted
representation (transport by a unitary equivalence
if necessary), of zero Jacquet norm. By
Theorem~\ref{thm:all-cell-relative-factorization}
that difference has zero orbital integral on every
relevant nonzero-Gram cell. This gives the exact
expansion
\[
 D_\ell(\gamma)\mathrm{OrbInt}(F_\eta,g,\nu)
 =\sum_{|I|=\ell}a_{\eta,I}^{(\ell)}
             \mathrm{OrbInt}(F_{\eta+\mathbf1_I},g,\nu)
                                                        \tag{AP5}
\]
with precisely the previous coefficients and the same
sum bound \(\sum_I|a_{\eta,I}^{(\ell)}|\le b^{C_bm^2}\),
\(m=2s\). There is no limiting passage between strata.
In particular, if the original input coefficient has
zero Jacquet norm, all its orbital evaluations vanish.

For the conjugated minor the tensor factor is the
dual exterior representation: on the unitary group
the complex conjugate of a standard coefficient
is a dual coefficient, using the invariant Hermitian
form. The admissible shifts are \(-\mathbf1_I\).
On the radial space multiplication by
\(\overline{D_\ell}\) is adjoint to multiplication by
\(D_\ell\), as is seen by integrating their product;
both products have finite Jacquet norms under the
retained reserve, or by the finite dual-Pieri
decomposition and (AP2). Hence the coefficient from
\(\eta\) to \(\eta-\mathbf1_I\) is the conjugate of
the already bounded coefficient in the other
direction. Its column sum is at most
\(\binom n\ell b^{C_bm^2}\le b^{(C_b+1)m^2}\).
Use (AP3) exactly as above to transport this equality
to the nonzero-Gram cells. This supplies the
all-cell equality in the original
\(\mathrm{OneStepCellExpansion}\) with an enlargement
of its uniform coefficient constant by one, under the
already established regular expansion hypotheses.
The local pole estimate itself is not proved by
this transfer.

\subsection{Weighted rotated summation with a rank-dependent pole exponent}
\label{sec:rank-dependent-rotated-sum}

The source's intermediate \(\mathrm{WeightedExpansionRotated}\)
uses a shift radius and an inverse-Gram exponent chosen
before the rank. Neither restriction is needed for the
actual \(\mathrm{JointNormBoundRotatedUniform}\) consumer.
The following argument replaces that intermediate
interface. It uses positive minors only.

\begin{theorem}[Reduction to the near-parameter unweighted orbital bound]
\label{thm:rotated-sum-from-near-bp}
Let \(a_w\ge0\), and retain the original
\(\mathrm{BData}\) at distance zero, with
\(\kappa=kaRe\), the original section and compact data.
Retain the actual range with nonempty cloud-gap set
(in particular the standing interlacing range \(k\ge3\)).
Suppose the unweighted near-parameter pointwise
orbital estimate holds with one constant \(C_*\ge1\),
on the same data and their admissible shifted family:
\[
 n_v(\gamma)\operatorname{orbSq}(F_\eta,
                  \mathrm{turnPt}(\mathrm{ptLoc}(\gamma)))
 \le b^{C_*m^4}
   \max\{1,G_{\rm loc}^{-1}\}^{C_*m^2}jNorm2(F_\eta).
                                                        \tag{AP6}
\]
Here \(G_{\rm loc}\) is the original positive Gram
polynomial at a represented point; it is exactly the
\(Gpt\) in \(\mathrm{BpIneq}\). The estimate is required
when the base parameter satisfies its original standard
bounds and all gaps at least \(200\|\eta\|_1\).
The required family consists only of the final coefficients in
the positive-minor expansion below, whose shifts satisfy
\(\|\eta\|_1\le sP\). In particular the prescribed character
scale, or its rigorously transported estimate, is retained.
Then the original weighted rotated point sum is bounded by
\[
 BSumR\bigl(a\mapsto normWt(a)^{a_ws^2}\bigr)
             \le b^{C_{\rm rot}m^5 f}\,jSqG.            \tag{AP7}
\]
The constants and strengthened base all-gap threshold
depend only on \(a_w,C_*,C_h\) and the already established
Pieri constants, not on \(r,k,f\). This is a conditional
reduction of the actual weighted-sum supplier; it does
not assume the stronger fixed-radius weighted expansion.
\end{theorem}

\begin{proof}
Put \(\rho=\operatorname{rd}(E)\), \(p=a_ws^2\), and
\[
       A_s=\lceil C_*m^2\rceil,\qquad L_s=\lceil p\rceil,
       \qquad M=\lceil 2a_w+4C_*+8\rceil,\qquad P=M s^3 .
\]
The original size hypothesis gives \(b\ge2\) and
\(n\le b\), and \(kaRe=2\pi/(g\rho)\) with
\(1\le\rho\), \(g\ge1\). From \((kaRe)^{-1}\le b^{C_h}\)
one gets
\(\rho\le b^{C_\rho}\), where
\(C_\rho=\max(0,C_h)+3\).
This is the original character scale, not a new
bound on the finite level.

Fix an admitted point \(\gamma\) and place \(v\).
If there is no relevant represented orbital datum,
its orbital supremum is zero; if the original
coefficient has zero Jacquet norm it vanishes by
(AP3). Otherwise normalize the input to radial
norm one. The argument below can then be multiplied
by its original squared Jacquet norm. Denote by
\(\mathcal E\) the block ends, including \(s\), and by
\(d_\ell=|\sigma_v(D_\ell^{\rm glob})|\) the
unnormalized minor moduli. Set \(u_\ell\) equal to
the corresponding minor moduli on the rotated
local point. The source's local scale and unit
rotation give exactly
\[
 d_\ell=\rho^{\ell(2r+1-\ell)}u_\ell,\qquad
 G_{\rm loc}=\rho^{-2rs}G_v,\qquad
 G_v=|\sigma_v(G_{\rm glob})| .
\]
The absolute value in the last equality agrees
with the positive sign at a represented nonzero-Gram
point. The relation for \(G_{\rm loc}\) follows
because the leading minors of lengths \(r,s\)
have the same scale exponent \(rs\).
In particular \(0<G_v\le d_s^2\), and
\(d_\ell\le\rho^{s^2}u_\ell\).
The source weights are
\[
 w_v=\prod_{\ell\in\mathcal E}(1+d_\ell^2)^2,\quad
 dLoc_v=d_s^2,\quad W_v=G_v\prod_{\ell\in\mathcal E}d_\ell^2 .
\]
Choose \(q_\ell=L_s+4\) for each block end, with
an additional \(A_s+L_s\) at \(s\).
For a fixed point choose
\(\alpha_\ell=q_\ell\) when \(u_\ell\ge1\)
and zero otherwise; non-end exponents are zero.
Since \(G_v\le d_s^2\), elementary inequalities give
\[
 w_v(1+d_s^2)^{A_s/2}W_v^{p/2}
       \le 2^{2s+A_s/2}\rho^{s^2P}
                    \prod_\ell u_\ell^{\alpha_\ell}.    \tag{AP8}
\]
Indeed \((1+d^2)^2\le4\max(1,d)^4\),
\((1+d_s^2)^{A_s/2}\le2^{A_s/2}\max(1,d_s)^{A_s}\),
and the two \(p\) factors at the last minor account
for the extra \(L_s\). The total monomial degree is
\[
 \sum_\ell\alpha_\ell
 \le s(L_s+4)+A_s+L_s
 \le (2a_w+4C_*+7)s^3\le P.
\]
Thus the monomial is allowed by the actual higher-minor
Pieri theorem.

All the required shifted parameters lie within
\(sP=M s^4\) of the base. Let \(C_g^{\rm mom}\)
be the fixed threshold constant for the existing
moment theorem, and \(C_q\ge200\) its constant in
\(q_0\). A base all-gap lower bound \(C_gm^4\)
with \(C_g\) sufficiently large in terms of these
constants and \(M\) ensures both
\[
 C_gm^4\ge C_g^{\rm mom} P m,\qquad
 C_gm^4\ge200sP.
\]
For the cloud gap (a nonempty finite set) the
same all-gap bound applies. The original
noncompact-gap bound gives
\[
 \frac{h_{NC}}{C_q\sqrt m}
 \ge \frac{C_g m^4(3sn^2+n^3)}{C_q}.
\]
Taking \(C_g\) also at least \(C_q^2(M+2)\)
therefore gives the full reserve
\(q_0\ge s(P+s+1)\). This is precisely the reserve
needed for the iterated positive-minor theorem,
including the one-step margin at each intermediate
shift. All the earlier representation and measure
parameters remain fixed except the discrete-series
parameter and its chosen coefficient.

Iterate (AP5) along the selected monomial, with
the original path-sum bound from the higher
Jacquet norm theorem. The resulting orbital
linear combination has coefficient sum at most
\(b^{C_bm^2P}\); every nonzero final coefficient
has normalized radial norm one. Take norms on
each actual orbital datum, use the triangle
inequality, and take the supremum in the original
\(\operatorname{orbSq}\). Using (AP6) for the
shifted coefficients and then (AP8) yields
\[
 \begin{split}
 &w_v^2(1+d_s^2)^{A_s}W_v^p\,
           n_v\,\operatorname{orbSq}(F)\\
 &\quad\le
   2^{4s+A_s}\rho^{2s^2P}
       b^{2C_bm^2P+C_*m^4}
       \max(1,\rho^{2rs}/G_v)^{C_*m^2}jNorm2(F)\\
 &\quad\le b^{C_2m^5}\max(1,1/G_v)^{A_s}jNorm2(F).
 \end{split}
 \tag{AP9}
\]
Here \(C_2\) is independent of the arithmetic datum.
For example one may take
\[
 C_2=10+(2C_\rho+8C_b)M
                  +(8C_\rho+20)C_* .
\]
To check this bound, use \(P=Ms^3\),
\(A_s\le(4C_*+1)s^2\), \(m=2s\), \(2\le b\),
and bound the extra inverse-Gram scaling by
\(\rho^{2rs C_*m^2}\le b^{8C_\rho C_*s^4}\).
All powers of \(s\) in the remaining exponent
are at most \(s^5\). Thus (AP9) has exactly the
one-place form needed for the source's cancellation,
but with a rank-dependent \(A_s\).

We reproduce that cancellation with no uniform
restriction on \(A_s\). Since \(G_{\rm glob}\) is
a nonzero integral element and lies in the real
subfield, the product of its absolute values over
the CM places is at least one. Therefore
\[
 \prod_v\max(1,G_v^{-1})^{A_s}
 \le\prod_v\max(1,G_v)^{A_s}
 \le\prod_v(1+dLoc_v)^{A_s}.
\]
Multiply (AP9) over the \(f\) places and cancel the
positive finite product on the right. It follows
that the summand in the original weighted \(BSumR\)
is bounded by
\[
    wPt(\gamma)^{-1} b^{C_2m^5 f}jSqG.                 \tag{AP10}
\]
This argument applies to extended nonnegative
quantities as well: the radial norms are finite,
a zero norm was dealt with by (AP3), and all
canceled scalar factors are finite and positive.

For completeness there is a direct elementary
uniform bound on the remaining reciprocal sum.
For each algebraic integer \(x\in\mathcal O_E\)
put \(H(x)=\prod_v(1+|\sigma_v(x)|^2)^{-2}\).
Partition each complex coordinate into half-unit
squares with sides parallel to the axes. There
is at most one embedded algebraic integer per
product of squares: a nonzero difference would
have all its complex moduli less than one, in
contradiction to its nonzero integral norm.
The sum of the suprema of the one-coordinate
weight over these squares is less than \(256\).
Indeed label squares by \((i,j)\in\mathbb Z^2\)
and put \(h=\max(|i|,|j|)\); the \(8h\) squares
at height \(h\ge4\) have moduli at least \(h/4\).
Their total weight is at most
\(2048\sum_{h\ge4}h^{-3}\le2048/18\), and the
\(49\) remaining squares each contribute at most
one. Thus \(\sum_x H(x)\le256^f\).

For a fixed composition of \(s\) into block
sizes \(t_i\), the tuple of nonzero block-end
minors consists of algebraic integers, by
invariance of leading minors under the integral
cell condition. It determines the block scalar
at step \(i\) up to at most \(t_i\) possibilities:
its \(t_i\)-th power is the ratio of the two
successive minors up to sign. Consequently the
fiber has cardinality at most
\(\prod_i t_i\le2^{s-1}\). There are \(2^{s-1}\)
compositions, and at most \(s\) block ends.
Using the preceding bound independently for
these integral minors proves
\[
        \sum_{\gamma\ {\rm admitted}}wPt(\gamma)^{-1}
                 \le2^{2(s-1)}256^{fs}
                 \le b^{10fs}.
\]
No lower bound for an individual complex embedding
has been assumed. Combining this with (AP10)
proves (AP7), for example with
\(C_{\rm rot}=C_2+10\).
\end{proof}

The pole exponent \(A_s\) and the shift radius \(sP\)
were chosen as functions of the rank, but the final
constant and the lower bound on \(C_g\) were chosen
before the datum. This is the quantifier order of the
weighted-sum consumer. It bypasses the extra fixed
\(A\) and fixed \(\mathrm{dist}'\) requirements of the
intermediate \(\mathrm{WeightedExpansionRotated}\).
The unweighted estimate (AP6) for precisely this finite shifted
family is supplied by
Theorem~\ref{thm:archimedean-pointwise-sufficient} and
Proposition~\ref{prop:archimedean-ap7-reserve}. The latter
checks its explicit spectral reserve using the consumer's
original fourth-power all-gap hypothesis. The exact amplitude
and row-norm comparison needed for
\(\mathrm{StepThmBPinnedG}\) is supplied in
Section~\ref{sec:flag-exact-norm-comparison}.
